% Part: first-order-logic
% Chapter: beyond
% Section: many-sorted-logic

\documentclass[../../../include/open-logic-section]{subfiles}

\begin{document}

\olfileid{fol}{byd}{msl}

\olsection{বহুজাতীয় যুক্তিবিদ্যা}

প্রথম-ক্রমের যুক্তিবিদ্যায় চলরাশি ও পরিমাণসূচক একটি মাত্র
!!{domain}-এর উপর বিস্তৃত হয়। কিন্তু অনেক সময় একাধিক (পরস্পর
বিচ্ছিন্ন) !!{domain}s রাখা উপকারী: যেমন, সংখ্যার !!a{domain},
জ্যামিতিক বস্তুর !!a{domain}, সংখ্যা থেকে সংখ্যায় অপেক্ষকগুলির
!!a{domain}, আবেলীয় গোষ্ঠীগুলির !!a{domain}, ইত্যাদি রাখতে চাইতে পারি।

বহুজাতীয় যুক্তিবিদ্যা এই ধরনের কাঠামো দেয়। শুরুতে “জাতি”-র একটি
তালিকা নেওয়া হয়---কোনো বস্তুর “জাতি” বলে দেয়, বস্তুটির কোন
“!!{domain}”-এ থাকার কথা। তারপর প্রত্যেক জাতির জন্য আলাদা
!!{variable}s ও পরিমাণসূচক এবং (সাধারণত) প্রত্যেক জাতির জন্য আলাদা
!!{identity} থাকে। অপেক্ষক ও সম্পর্কও তাদের আর্গুমেন্ট হিসেবে নেওয়া
বস্তুর জাতি অনুসারে “টাইপযুক্ত” হয়। অন্য সব ক্ষেত্রে প্রথম-ক্রমের
যুক্তিবিদ্যার পরিচিত বিধিগুলিই রাখা হয়; শুধু পরিমাণসূচক-বিধিগুলির
একটি করে সংস্করণ প্রত্যেক জাতির জন্য পুনরাবৃত্ত হয়।

উদাহরণ হিসেবে আন্তর্জাতিক সম্পর্ক অধ্যয়নের জন্য আমরা দুই জাতির
বস্তুবিশিষ্ট একটি ভাষা বেছে নিতে পারি: ফরাসি নাগরিক ও জার্মান নাগরিক।
প্রথম জাতির বস্তুর জন্য “মদ্যপান করে” নামে একটি এক-স্থানীয় সম্পর্ক,
দ্বিতীয় জাতির বস্তুর জন্য “ভুর্স্ট খায়” নামে আরেকটি এক-স্থানীয়
সম্পর্ক, এবং “বহুজাতিক বিবাহিত দম্পতি গঠন করে” নামে একটি দ্বি-স্থানীয়
সম্পর্ক থাকতে পারে; শেষের সম্পর্কটির প্রথম আর্গুমেন্ট প্রথম জাতির এবং
দ্বিতীয় আর্গুমেন্ট দ্বিতীয় জাতির। ফরাসি নাগরিকদের উপর বিস্তৃত হতে
$a$, $b$, $c$ এবং জার্মান নাগরিকদের উপর বিস্তৃত হতে $x$, $y$, $z$
চলরাশি ব্যবহার করলে
\[
\lforall[a][\lforall[x][(\Atom{\Obj{MarriedTo}}{a,x} \lif
(\Atom{\Obj{DrinksWine}}{a} \lor \lnot \Atom{\Obj{EatsWurst}}{x})]]
\]
% BN-SRC-109 TARGET-CORRECTION-BEGIN
\par\smallskip\noindent\textnormal{\small\textbf{উৎস-সংশোধন (BN-SRC-109):} হিমায়িত ইংরেজি উৎসে অন্তঃস্থ সার্বিক পরিমাণসূচকের চলরাশি-আর্গুমেন্টের উদ্বোধনী বর্গবন্ধনী অনুপস্থিত ছিল: \texttt{\string\lforall x]}। বাংলা সূত্রে \texttt{\string\lforall[x]} পুনরুদ্ধার করা হয়েছে।} % BN-SRC-109 TARGET-CORRECTION-END
দাবি করে: কোনো ফরাসি ব্যক্তি কোনো জার্মানকে বিয়ে করে থাকলে হয় সেই
ফরাসি ব্যক্তি মদ্যপান করেন, নয় সেই জার্মান ব্যক্তি ভুর্স্ট খান না।

বহুজাতীয় যুক্তিবিদ্যাকে স্বাভাবিক উপায়ে প্রথম-ক্রমের যুক্তিবিদ্যার
মধ্যে নিবেশিত করা যায়। বহুজাতীয় !!{domain}s-গুলির সব বস্তুকে একত্র
করে একটি প্রথম-ক্রমের !!{domain} নেওয়া হয়, এক-স্থানীয়
!!{predicate}s দিয়ে জাতিগুলির হিসাব রাখা হয়, এবং পরিমাণসূচকগুলিকে
আপেক্ষিক করা হয়। যেমন, উপরের উদাহরণের অনুরূপ প্রথম-ক্রমের ভাষায়
বর্ণিত অন্য সম্পর্কগুলির সঙ্গে “$\Obj{German}$” ও
“$\Obj{French}$” নামে এক-স্থানীয় !!{predicate}s থাকবে, আর জাতি-সংক্রান্ত
শর্তগুলি মুছে যাবে। জার্মান জাতির !!a{variable} $x$-এর জন্য
বহুজাতীয় পরিমাণসূচক $\lforall[x][!A]$ অনূদিত হবে
\[
\lforall[x][(\Atom{\Obj{German}}{x} \lif !A)].
\]
জাতিগুলি যে পৃথক, তা নিশ্চিত করার স্বতঃসিদ্ধ---যেমন,
$\lforall[x][\lnot (\Atom{\Obj{German}}{x} \land
  \Atom{\Obj{French}}{x})]$---যোগ করতে হবে; একই সঙ্গে এমন স্বতঃসিদ্ধও
লাগবে যা নিশ্চিত করে যে “মদ্যপান করে” কেবল
$\Atom{\Obj{French}}{x}$ বিধেয়টি পরিতৃপ্ত করা বস্তুর ক্ষেত্রেই সত্য,
ইত্যাদি। এই প্রথা ও স্বতঃসিদ্ধগুলি মানলে বহুজাতীয় !!{sentence}s যে
প্রথম-ক্রমের !!{sentence}s-এ অনূদিত হয় এবং বহুজাতীয়
!!{derivation}s যে প্রথম-ক্রমের !!{derivation}s-এ অনূদিত হয়, তা
দেখানো কঠিন নয়। আবার বহুজাতীয় !!{structure}s-ও অনুরূপ
প্রথম-ক্রমের !!{structure}s-এ “অনূদিত” হয় এবং বিপরীত অনুবাদও চলে;
তাই বহুজাতীয় যুক্তিবিদ্যার জন্যও একটি সম্পূর্ণতা উপপাদ্য পাওয়া যায়।

\end{document}
